\documentclass[11pt,a4paper]{article}
\usepackage[utf8]{inputenc}
\usepackage[T1]{fontenc}
\usepackage[spanish,es-minimal,es-noshorthands]{babel}
\AtBeginDocument{\renewcommand{\sin}{\operatorname{sen}}\renewcommand{\arcsin}{\operatorname{arcsen}}}
\usepackage{lmodern}
\usepackage[sfdefault,lining]{FiraSans}
\usepackage[a4paper,top=2.3cm,bottom=2.2cm,left=1.8cm,right=1.8cm,headheight=1.3cm,headsep=0.4cm,footskip=1cm]{geometry}
\usepackage{amsmath,amssymb}
\usepackage{icomma}
\usepackage{graphicx}
\usepackage[table]{xcolor}
\usepackage{tikz}
\usepackage{fontawesome5}
\usepackage{ccicons}
\usepackage{fancyhdr}
\usepackage{lastpage}
\usepackage{tasks}
\usepackage{needspace}
\usepackage[most]{tcolorbox}
\usepackage{array}
\usepackage[hidelinks]{hyperref}
\AtBeginDocument{\hypersetup{pdftitle={Ficha de ejercicios de la unidad 3: Introducción al álgebra. Polinomios. 2º ESO},pdfsubject={Matemáticas 2º ESO},pdfkeywords={lenguaje algebraico, valor numérico, monomios, monomios semejantes, operaciones con monomios, potencia de un monomio, polinomios, suma y resta de polinomios, producto de polinomios, identidades notables, factor común}}}

\definecolor{navy}{HTML}{1B2748}
\definecolor{accent}{HTML}{E8603A}
\definecolor{teal}{HTML}{1F9E90}
\definecolor{soft}{HTML}{EEF1F8}
\definecolor{muted}{HTML}{5D6685}

\setlength{\parindent}{0pt}
\setlength{\parskip}{3pt}
\settasks{label=\alph*), label-format=\bfseries\color{navy}, label-width=1.3em, item-indent=1.8em, column-sep=0.6em, after-item-skip=5pt, before-skip=3pt, after-skip=4pt}

% ---------- capçalera i peu ----------
\newcommand{\llicencia}{\href{https://creativecommons.org/licenses/by-nc-sa/4.0/deed.es}{\ccbyncsa}\ \ Licencia Creative Commons Atribución-NoComercial-CompartirIgual 4.0 Internacional (CC BY-NC-SA 4.0)}
\pagestyle{fancy}
\fancyhf{}
\fancyhead[L]{\raisebox{-0.35\height}{\includegraphics[height=0.95cm]{logo-gva.png}}}
\fancyhead[C]{\small\color{muted}Ficha de ejercicios · Unidad 3: Introducción al álgebra. Polinomios · 2\textsuperscript{o} ESO}
\fancyhead[R]{\raisebox{-0.35\height}{\includegraphics[height=0.95cm]{logo-vermellar.png}}}
\fancyfoot[L]{\scriptsize\color{muted}\llicencia}
\fancyfoot[R]{\small\color{muted}\thepage}
\renewcommand{\headrulewidth}{0.4pt}
\renewcommand{\headrule}{\color{muted!40}\hrule height \headrulewidth}
\fancypagestyle{primera}{\fancyhf{}\renewcommand{\headrulewidth}{0pt}
  \fancyfoot[L]{\scriptsize\color{muted}\llicencia}
  \fancyfoot[R]{\small\color{muted}\thepage}}

% ---------- nivells i exercicis ----------
\newcommand{\estrelles}[1]{{\color{accent}\ifcase#1\or\faStar\faStar[regular]\faStar[regular]\or\faStar\faStar\faStar[regular]\or\faStar\faStar\faStar\fi}}
\newcounter{exer}
\newcommand{\ex}[1]{\par\needspace{4\baselineskip}\medskip\stepcounter{exer}\noindent{\color{navy}\bfseries\theexer.}\ \ {\small\estrelles{#1}}\ \ }

\newcommand{\seccio}[2]{\par\needspace{10\baselineskip}\bigskip
  \begin{tcolorbox}[colback=navy,colframe=navy,coltext=white,arc=2mm,boxrule=0pt,left=3mm,right=3mm,top=1.2mm,bottom=1.2mm]
  {\large\bfseries #1\quad #2}
  \end{tcolorbox}\vspace{-2pt}}

\newcolumntype{C}[1]{>{\centering\arraybackslash}p{#1}}
% quadre de recordatori (només a la fitxa d'exercicis)
\newcommand{\recorda}[1]{\par\smallskip\begin{tcolorbox}[colback=accent!7,colframe=accent!55,arc=2mm,boxrule=0.5pt,left=2.5mm,right=2.5mm,top=1mm,bottom=1mm]\small{\color{accent}\faLightbulb[regular]}\ \textbf{Recuerda:}\ #1\end{tcolorbox}}
\newcommand{\res}[1]{#1}\newcommand{\rest}[1]{#1}

\def\fmneed{12}
% =====================================================================
% Capçalera FentMates de les fitxes: un color per a cada curs.
% Els colors són els de la pàgina inicial (LEVELS a assets/js/materials.js);
% el logo ve de _eines/logo/pdf/ (python3 _eines/logo/genera.py).
%
% Ús, al final del preàmbul de la unitat:
%   % =====================================================================
% Capçalera FentMates de les fitxes: un color per a cada curs.
% Els colors són els de la pàgina inicial (LEVELS a assets/js/materials.js);
% el logo ve de _eines/logo/pdf/ (python3 _eines/logo/genera.py).
%
% Ús, al final del preàmbul de la unitat:
%   % =====================================================================
% Capçalera FentMates de les fitxes: un color per a cada curs.
% Els colors són els de la pàgina inicial (LEVELS a assets/js/materials.js);
% el logo ve de _eines/logo/pdf/ (python3 _eines/logo/genera.py).
%
% Ús, al final del preàmbul de la unitat:
%   \input{../comu/fentmates.tex}
%   \fmcurs{1eso}
% i, al començament del document:
%   \capcalera{Matemàtiques}{1r ESO}{Fitxa d'exercicis · Unitat 2}{Divisibilitat}
% Opcional abans de \input: \def\fmneed{14} (espai mínim abans de cada secció).
% =====================================================================
\makeatletter
\def\fm@logodir{../../logo/pdf}
\providecommand{\fmneed}{10}
\def\fm@cursnom{}

% clau del curs, color i color del text damunt del color
\newcommand{\fm@curs}[3]{\@namedef{fm@color@#1}{#2}\@namedef{fm@text@#1}{#3}}
\fm@curs{1eso}{E6B84A}{navy}
\fm@curs{2eso}{3B82F6}{white}
\fm@curs{3eso}{FF7A50}{white}
\fm@curs{4eso-a}{A78BFA}{white}
\fm@curs{4eso-b}{6D28D9}{white}
\fm@curs{1batx-mi}{9B1C1C}{white}
\fm@curs{1batx-ccss}{EF6B6B}{white}
\fm@curs{2batx-mii}{1E6B47}{white}
\fm@curs{2batx-ccss}{2BB3A3}{white}

\newcommand{\fmcurs}[1]{%
  \@ifundefined{fm@color@#1}{\PackageError{fentmates}{Curs desconegut: #1}{}}{}%
  \edef\fm@tmp{\noexpand\definecolor{curs}{HTML}{\@nameuse{fm@color@#1}}%
               \noexpand\colorlet{curstext}{\@nameuse{fm@text@#1}}}\fm@tmp
  \colorlet{cursfosc}{curs!72!black}%
  \def\fm@logo{\fm@logodir/vertical-#1.pdf}%
  \def\fm@marca{\fm@logodir/marca-#1.pdf}}

% ---------- capçalera de la primera pàgina ----------
% #1 matèria (petit), #2 curs, #3 tipus de document i unitat, #4 títol
\newcommand{\capcalera}[4]{%
  \gdef\fm@cursnom{#2}%
  \par\noindent
  \begin{tcolorbox}[colback=curs,colframe=curs,boxrule=0pt,arc=2.2mm,width=3.5cm,height=2.15cm,
      valign=center,left=3mm,right=2mm,top=1mm,bottom=1mm,nobeforeafter,box align=center]
    \color{curstext}{\footnotesize #1}\par\vspace{1pt}{\fontsize{19}{21}\selectfont\bfseries #2\par}%
  \end{tcolorbox}\hspace{5mm}%
  \begin{minipage}[c]{\dimexpr\linewidth-3.5cm-5mm-2.9cm\relax}
    {\color{cursfosc}\small\bfseries\MakeUppercase{#3}}\par\vspace{2pt}
    {\color{navy}\fontsize{20}{23}\selectfont\bfseries #4\par}
  \end{minipage}\hfill
  \begin{minipage}[c]{2.6cm}\raggedleft\includegraphics[width=2.6cm]{\fm@logo}\end{minipage}\par
  \vspace{2pt}}

% ---------- barres de secció del color del curs ----------
\renewcommand{\seccio}[2]{\par\needspace{\fmneed\baselineskip}\bigskip
  \begin{tcolorbox}[colback=curs,colframe=curs,coltext=curstext,arc=2mm,boxrule=0pt,left=3mm,right=3mm,top=1.2mm,bottom=1.2mm]
  {\large\bfseries #1\quad #2}
  \end{tcolorbox}\vspace{-2pt}}

% ---------- estrelles de dificultat del color del curs (a Batxillerat no se'n posen) ----------
\@ifundefined{estrelles}{}{%
  \renewcommand{\estrelles}[1]{{\color{curs}\ifcase#1\or\faStar\faStar[regular]\faStar[regular]\or\faStar\faStar\faStar[regular]\or\faStar\faStar\faStar\fi}}}

% ---------- capçalera i peu de les pàgines ----------
% primera pàgina: logos institucionals al peu
\fancypagestyle{primera}{\fancyhf{}\renewcommand{\headrulewidth}{0pt}
  \fancyfoot[L]{\raisebox{-0.32\height}{\includegraphics[height=0.78cm]{logo-gva.png}}\hspace{5mm}%
    \raisebox{-0.32\height}{\includegraphics[height=0.78cm]{logo-vermellar.png}}}
  \fancyfoot[R]{\scriptsize\color{muted}\href{https://creativecommons.org/licenses/by-nc-sa/4.0/deed.ca}{\ccbyncsa}\ \ CC BY-NC-SA 4.0 · fentmates.org\qquad\small\thepage}}
% la resta: el símbol a l'esquerra i el curs a la dreta, amb el títol al centre com abans
\fancyhead[L]{\raisebox{-0.3\height}{\includegraphics[height=0.62cm]{\fm@marca}}}
\fancyhead[R]{\small\bfseries\color{cursfosc}\fm@cursnom}
\renewcommand{\headrule}{\color{curs!45}\hrule height \headrulewidth}
\makeatother

%   \fmcurs{1eso}
% i, al començament del document:
%   \capcalera{Matemàtiques}{1r ESO}{Fitxa d'exercicis · Unitat 2}{Divisibilitat}
% Opcional abans de \input: \def\fmneed{14} (espai mínim abans de cada secció).
% =====================================================================
\makeatletter
\def\fm@logodir{../../logo/pdf}
\providecommand{\fmneed}{10}
\def\fm@cursnom{}

% clau del curs, color i color del text damunt del color
\newcommand{\fm@curs}[3]{\@namedef{fm@color@#1}{#2}\@namedef{fm@text@#1}{#3}}
\fm@curs{1eso}{E6B84A}{navy}
\fm@curs{2eso}{3B82F6}{white}
\fm@curs{3eso}{FF7A50}{white}
\fm@curs{4eso-a}{A78BFA}{white}
\fm@curs{4eso-b}{6D28D9}{white}
\fm@curs{1batx-mi}{9B1C1C}{white}
\fm@curs{1batx-ccss}{EF6B6B}{white}
\fm@curs{2batx-mii}{1E6B47}{white}
\fm@curs{2batx-ccss}{2BB3A3}{white}

\newcommand{\fmcurs}[1]{%
  \@ifundefined{fm@color@#1}{\PackageError{fentmates}{Curs desconegut: #1}{}}{}%
  \edef\fm@tmp{\noexpand\definecolor{curs}{HTML}{\@nameuse{fm@color@#1}}%
               \noexpand\colorlet{curstext}{\@nameuse{fm@text@#1}}}\fm@tmp
  \colorlet{cursfosc}{curs!72!black}%
  \def\fm@logo{\fm@logodir/vertical-#1.pdf}%
  \def\fm@marca{\fm@logodir/marca-#1.pdf}}

% ---------- capçalera de la primera pàgina ----------
% #1 matèria (petit), #2 curs, #3 tipus de document i unitat, #4 títol
\newcommand{\capcalera}[4]{%
  \gdef\fm@cursnom{#2}%
  \par\noindent
  \begin{tcolorbox}[colback=curs,colframe=curs,boxrule=0pt,arc=2.2mm,width=3.5cm,height=2.15cm,
      valign=center,left=3mm,right=2mm,top=1mm,bottom=1mm,nobeforeafter,box align=center]
    \color{curstext}{\footnotesize #1}\par\vspace{1pt}{\fontsize{19}{21}\selectfont\bfseries #2\par}%
  \end{tcolorbox}\hspace{5mm}%
  \begin{minipage}[c]{\dimexpr\linewidth-3.5cm-5mm-2.9cm\relax}
    {\color{cursfosc}\small\bfseries\MakeUppercase{#3}}\par\vspace{2pt}
    {\color{navy}\fontsize{20}{23}\selectfont\bfseries #4\par}
  \end{minipage}\hfill
  \begin{minipage}[c]{2.6cm}\raggedleft\includegraphics[width=2.6cm]{\fm@logo}\end{minipage}\par
  \vspace{2pt}}

% ---------- barres de secció del color del curs ----------
\renewcommand{\seccio}[2]{\par\needspace{\fmneed\baselineskip}\bigskip
  \begin{tcolorbox}[colback=curs,colframe=curs,coltext=curstext,arc=2mm,boxrule=0pt,left=3mm,right=3mm,top=1.2mm,bottom=1.2mm]
  {\large\bfseries #1\quad #2}
  \end{tcolorbox}\vspace{-2pt}}

% ---------- estrelles de dificultat del color del curs (a Batxillerat no se'n posen) ----------
\@ifundefined{estrelles}{}{%
  \renewcommand{\estrelles}[1]{{\color{curs}\ifcase#1\or\faStar\faStar[regular]\faStar[regular]\or\faStar\faStar\faStar[regular]\or\faStar\faStar\faStar\fi}}}

% ---------- capçalera i peu de les pàgines ----------
% primera pàgina: logos institucionals al peu
\fancypagestyle{primera}{\fancyhf{}\renewcommand{\headrulewidth}{0pt}
  \fancyfoot[L]{\raisebox{-0.32\height}{\includegraphics[height=0.78cm]{logo-gva.png}}\hspace{5mm}%
    \raisebox{-0.32\height}{\includegraphics[height=0.78cm]{logo-vermellar.png}}}
  \fancyfoot[R]{\scriptsize\color{muted}\href{https://creativecommons.org/licenses/by-nc-sa/4.0/deed.ca}{\ccbyncsa}\ \ CC BY-NC-SA 4.0 · fentmates.org\qquad\small\thepage}}
% la resta: el símbol a l'esquerra i el curs a la dreta, amb el títol al centre com abans
\fancyhead[L]{\raisebox{-0.3\height}{\includegraphics[height=0.62cm]{\fm@marca}}}
\fancyhead[R]{\small\bfseries\color{cursfosc}\fm@cursnom}
\renewcommand{\headrule}{\color{curs!45}\hrule height \headrulewidth}
\makeatother

%   \fmcurs{1eso}
% i, al començament del document:
%   \capcalera{Matemàtiques}{1r ESO}{Fitxa d'exercicis · Unitat 2}{Divisibilitat}
% Opcional abans de \input: \def\fmneed{14} (espai mínim abans de cada secció).
% =====================================================================
\makeatletter
\def\fm@logodir{../../logo/pdf}
\providecommand{\fmneed}{10}
\def\fm@cursnom{}

% clau del curs, color i color del text damunt del color
\newcommand{\fm@curs}[3]{\@namedef{fm@color@#1}{#2}\@namedef{fm@text@#1}{#3}}
\fm@curs{1eso}{E6B84A}{navy}
\fm@curs{2eso}{3B82F6}{white}
\fm@curs{3eso}{FF7A50}{white}
\fm@curs{4eso-a}{A78BFA}{white}
\fm@curs{4eso-b}{6D28D9}{white}
\fm@curs{1batx-mi}{9B1C1C}{white}
\fm@curs{1batx-ccss}{EF6B6B}{white}
\fm@curs{2batx-mii}{1E6B47}{white}
\fm@curs{2batx-ccss}{2BB3A3}{white}

\newcommand{\fmcurs}[1]{%
  \@ifundefined{fm@color@#1}{\PackageError{fentmates}{Curs desconegut: #1}{}}{}%
  \edef\fm@tmp{\noexpand\definecolor{curs}{HTML}{\@nameuse{fm@color@#1}}%
               \noexpand\colorlet{curstext}{\@nameuse{fm@text@#1}}}\fm@tmp
  \colorlet{cursfosc}{curs!72!black}%
  \def\fm@logo{\fm@logodir/vertical-#1.pdf}%
  \def\fm@marca{\fm@logodir/marca-#1.pdf}}

% ---------- capçalera de la primera pàgina ----------
% #1 matèria (petit), #2 curs, #3 tipus de document i unitat, #4 títol
\newcommand{\capcalera}[4]{%
  \gdef\fm@cursnom{#2}%
  \par\noindent
  \begin{tcolorbox}[colback=curs,colframe=curs,boxrule=0pt,arc=2.2mm,width=3.5cm,height=2.15cm,
      valign=center,left=3mm,right=2mm,top=1mm,bottom=1mm,nobeforeafter,box align=center]
    \color{curstext}{\footnotesize #1}\par\vspace{1pt}{\fontsize{19}{21}\selectfont\bfseries #2\par}%
  \end{tcolorbox}\hspace{5mm}%
  \begin{minipage}[c]{\dimexpr\linewidth-3.5cm-5mm-2.9cm\relax}
    {\color{cursfosc}\small\bfseries\MakeUppercase{#3}}\par\vspace{2pt}
    {\color{navy}\fontsize{20}{23}\selectfont\bfseries #4\par}
  \end{minipage}\hfill
  \begin{minipage}[c]{2.6cm}\raggedleft\includegraphics[width=2.6cm]{\fm@logo}\end{minipage}\par
  \vspace{2pt}}

% ---------- barres de secció del color del curs ----------
\renewcommand{\seccio}[2]{\par\needspace{\fmneed\baselineskip}\bigskip
  \begin{tcolorbox}[colback=curs,colframe=curs,coltext=curstext,arc=2mm,boxrule=0pt,left=3mm,right=3mm,top=1.2mm,bottom=1.2mm]
  {\large\bfseries #1\quad #2}
  \end{tcolorbox}\vspace{-2pt}}

% ---------- estrelles de dificultat del color del curs (a Batxillerat no se'n posen) ----------
\@ifundefined{estrelles}{}{%
  \renewcommand{\estrelles}[1]{{\color{curs}\ifcase#1\or\faStar\faStar[regular]\faStar[regular]\or\faStar\faStar\faStar[regular]\or\faStar\faStar\faStar\fi}}}

% ---------- capçalera i peu de les pàgines ----------
% primera pàgina: logos institucionals al peu
\fancypagestyle{primera}{\fancyhf{}\renewcommand{\headrulewidth}{0pt}
  \fancyfoot[L]{\raisebox{-0.32\height}{\includegraphics[height=0.78cm]{logo-gva.png}}\hspace{5mm}%
    \raisebox{-0.32\height}{\includegraphics[height=0.78cm]{logo-vermellar.png}}}
  \fancyfoot[R]{\scriptsize\color{muted}\href{https://creativecommons.org/licenses/by-nc-sa/4.0/deed.ca}{\ccbyncsa}\ \ CC BY-NC-SA 4.0 · fentmates.org\qquad\small\thepage}}
% la resta: el símbol a l'esquerra i el curs a la dreta, amb el títol al centre com abans
\fancyhead[L]{\raisebox{-0.3\height}{\includegraphics[height=0.62cm]{\fm@marca}}}
\fancyhead[R]{\small\bfseries\color{cursfosc}\fm@cursnom}
\renewcommand{\headrule}{\color{curs!45}\hrule height \headrulewidth}
\makeatother

\fmcurs{2eso}

\begin{document}
\thispagestyle{primera}

\capcalera{Matemáticas}{2\textsuperscript{o} ESO}{Ficha de ejercicios · Unidad 3}{Introducción al álgebra.\\ Polinomios}
\vspace{10pt}
Nombre y apellidos: \rule{6.6cm}{0.4pt}\hfill Grupo: \rule{1.6cm}{0.4pt}\hfill Fecha: \rule{2.2cm}{0.4pt}
\vspace{8pt}

\begin{tcolorbox}[colback=soft,colframe=soft,arc=2mm,left=3mm,right=3mm,top=1.5mm,bottom=1.5mm]
\small Ejercicios graduados por niveles de dificultad:\quad \estrelles{1}\ básico \qquad \estrelles{2}\ medio \qquad \estrelles{3}\ avanzado.\\[2pt] Resuélvelos en el cuaderno mostrando todos los pasos y sin calculadora. Da los resultados reducidos y ordenados. Las secciones siguen el orden de los apuntes de la unidad (3.1--3.8).
\end{tcolorbox}

% =====================================================
\seccio{1}{El álgebra: el lenguaje algebraico}

\ex{1} Escribe en lenguaje algebraico (usa $x$ para el número):\begin{tasks}[after-item-skip=3pt](2)
  \task El triple de un número.
  \task Un número menos 8.
  \task La mitad de un número.
  \task El cuadrado de un número más 1.
  \task El doble de un número menos 5.
  \task El doble de la diferencia de un número y 5.
  \task El siguiente de un número.
  \task La suma de un número y su cuadrado.
  \task El cubo de un número.
  \task La cuarta parte de un número menos 3.
\end{tasks}\recorda{el orden de las palabras importa: «el doble de un número menos 5» es $2x-5$, pero «el doble de la diferencia de un número y 5» es $2(x-5)$.}

\ex{1} Escribe con palabras qué expresa cada expresión algebraica:\begin{tasks}(4)
  \task $x+5$
  \task $4x$
  \task $x^2-3$
  \task $\dfrac{x}{3}$
  \task $2x+1$
  \task $(x+2)^2$
  \task $x-y$
  \task $10-x$
\end{tasks}

\ex{1} Escribe la expresión algebraica de: a) el perímetro de un cuadrado de lado $x$; b) el área de un rectángulo de base $x$ y altura $5$; c) la edad de una persona dentro de 7 años si ahora tiene $x$; d) el precio de $x$ kilos de manzanas a $1{,}80$ € el kilo.

\ex{1} Calcula el valor numérico de cada expresión para el valor de $x$ indicado:\begin{tasks}[after-item-skip=5pt](3)
  \task $3x+5$ \ para $x=2$
  \task $7-2x$ \ para $x=-3$
  \task $x^{2}-4x$ \ para $x=5$
  \task $2x^{2}+x-1$ \ para $x=-2$
  \task $-x^{2}+6$ \ para $x=-3$
  \task $x^{3}-x$ \ para $x=-2$
\end{tasks}\recorda{los números negativos se sustituyen entre paréntesis: si $x=-3$, \ $x^2=(-3)^2=9$ pero $-x^2=-(-3)^2=-9$.}

\ex{2} Calcula el valor numérico:\begin{tasks}[after-item-skip=5pt](2)
  \task $2a+3b$ \ para $a=4,\ b=-1$
  \task $a^{2}-b^{2}$ \ para $a=-3,\ b=2$
  \task $5xy-x$ \ para $x=2,\ y=-1$
  \task $x^{2}y+xy^{2}$ \ para $x=-1,\ y=3$
\end{tasks}

\ex{2} Una compañía de telefonía cobra una cuota de 9 € al mes más $0{,}05$ € por cada minuto de llamada. Escribe la expresión de lo que pagamos en un mes si hablamos $x$ minutos y calcula cuánto pagamos si hablamos 120 minutos. ¿Y si hablamos 3 horas?

\ex{2} Con palillos formamos una hilera de triángulos: la figura 1 tiene 3 palillos, la 2 tiene 5, la 3 tiene 7… ¿Cuántos tiene la figura 10? Escribe una expresión para los palillos de la figura $n$. ¿Qué figura tiene 51 palillos?

\ex{3} Escribe tres números consecutivos cualesquiera usando $x$ y demuestra que su suma siempre es múltiplo de 3. ¿Es cierto también para la suma de cuatro números consecutivos?


% =====================================================
\seccio{2}{Monomios: partes y valor numérico}

\ex{1} Completa la tabla:\par\smallskip{\centering\renewcommand{\arraystretch}{1.35}\begin{tabular}{|c|C{2.6cm}|C{2.6cm}|C{1.6cm}|}\hline Monomio & Coeficiente & Parte literal & Grado\\\hline $-7x^3$ & & & \\\hline $xy^2$ & & & \\\hline $\tfrac{2}{3}a^2b$ & & & \\\hline $-x$ & & & \\\hline $5x^2y^3z$ & & & \\\hline $12$ & & & \\\hline\end{tabular}\par}\recorda{el grado de un monomio es la \textbf{suma} de los exponentes de todas las letras; si no ves el coeficiente, es 1 (o $-1$).}

\ex{1} ¿Cuáles de estas expresiones son monomios? Explica por qué las otras no lo son.\begin{tasks}(4)
  \task $3x^4$
  \task $\dfrac{5}{x}$
  \task $x+y$
  \task $-\dfrac{x^2}{4}$
  \task $\sqrt{x}$
  \task $-2ab^3$
  \task $x^{-1}$
  \task $9$
\end{tasks}

\ex{1} Agrupa los monomios semejantes: \ $4x^2,\ \ -3xy,\ \ 2x,\ \ -x^2,\ \ \tfrac{1}{2}xy,\ \ 7x,\ \ 5x^2y,\ \ -x,\ \ \tfrac{3}{4}x^2,\ \ x^2y$.

\ex{2} Calcula el valor numérico de cada monomio:\begin{tasks}[after-item-skip=5pt](2)
  \task $-3x^{2}$ \ para $x=4$
  \task $2x^{3}$ \ para $x=-2$
  \task $-a^{2}b$ \ para $a=-3,\ b=2$
  \task $4xy^{2}$ \ para $x=\tfrac{1}{2},\ y=-3$
\end{tasks}

\ex{2} Escribe un monomio: a) semejante a $-6x^2y$ con coeficiente positivo; b) de grado 5 con dos letras y coeficiente $-\tfrac{1}{2}$; c) opuesto a $3ab^2$; d) de grado 0.

\ex{3} Halla $m$ y $n$ para que los monomios $5x^{m}y^{3}$ y $-2x^{4}y^{n+1}$ sean semejantes. ¿Qué grado tienen?

\ex{3} Escribe un monomio de grado 4, en las letras $x$ e $y$, con coeficiente 2, que valga 54 para $x=3$, $y=1$.


% =====================================================
\seccio{3}{Operaciones con monomios}

\ex{1} Reduce:\begin{tasks}[after-item-skip=6pt](3)
  \task $5x+3x-2x$
  \task $4x^{2}-9x^{2}$
  \task $-a+6a-8a$
  \task $3xy-xy+5xy$
  \task $x^{3}+x^{3}+x^{3}$
  \task $7ab^{2}-7ab^{2}$
\end{tasks}\recorda{solo se suman y restan monomios \textbf{semejantes}: sumamos los coeficientes y \textbf{dejamos la misma parte literal} ($5x^2+3x^2=8x^2$, nunca $8x^4$).}

\ex{1} Reduce cuando sea posible. Si no se puede, explica por qué:\begin{tasks}(4)
  \task $3x^{2}+5x$
  \task $4x-x^{2}+2x$
  \task $6ab-2a+ab$
  \task $2x^{3}-x^{2}+4x^{3}-x$
\end{tasks}

\ex{1} Multiplica:\begin{tasks}[after-item-skip=6pt](3)
  \task $3x^{2}\cdot 5x^{4}$
  \task $(-2a)\cdot (7a^{3})$
  \task $(-4xy)\cdot (-3x^{2}y)$
  \task $6x^{2}\cdot (\tfrac{1}{3})x$
  \task $x^{3}y\cdot (-xy^{4})$
  \task $2a\cdot 3b\cdot (-a^{2})$
\end{tasks}\recorda{multiplicamos los coeficientes y \textbf{sumamos} los exponentes de cada letra: $x^m\cdot x^n=x^{m+n}$.}

\ex{1} Divide:\begin{tasks}[after-item-skip=6pt](3)
  \task $12x^{5}:(4x^{2})$
  \task $(-20a^{4}):(5a)$
  \task $18x^{3}y^{2}:(-6xy)$
  \task $(-9x^{6}):(-9x^{6})$
  \task $7x^{4}y:(14x^{2})$
  \task $(-15a^{3}b^{4}):(3ab^{4})$
\end{tasks}\recorda{dividimos los coeficientes y \textbf{restamos} los exponentes: $x^m:x^n=x^{m-n}$; y $x^0=1$.}

\ex{2} Calcula las potencias:\begin{tasks}[after-item-skip=6pt](3)
  \task $(2x^{3})^{2}$
  \task $(-3a^{2})^{3}$
  \task $(x^{2}y^{3})^{4}$
  \task $(-2xy^{2})^{4}$
  \task $(-a^{3}b)^{5}$
  \task $(\tfrac{1}{2}x^{2})^{3}$
\end{tasks}\recorda{elevamos el coeficiente y \textbf{multiplicamos} los exponentes: $(-2x^3y)^2=(-2)^2x^{6}y^{2}=4x^6y^2$. Cuidado con el signo si el exponente es impar.}

\ex{2} Reduce agrupando los términos semejantes:\begin{tasks}[after-item-skip=6pt](2)
  \task $5x^{2}-3x+2-x^{2}+7x-6$
  \task $-2a+4b-a-3b+5a$
  \task $x^{3}-2x+x^{2}-4x^{3}+3x+x^{2}$
  \task $3xy-x^{2}+2xy+4x^{2}-y$
\end{tasks}

\ex{2} ¿Cuáles de estos cocientes son monomios? Calcúlalos.\begin{tasks}(4)
  \task $8x^{4}:(2x)$
  \task $6x^{2}:(3x^{5})$
  \task $10a^{3}b:(5ab)$
  \task $4xy:(2y^{2})$
\end{tasks}

\ex{2} Opera y reduce (respeta la jerarquía de las operaciones):\begin{tasks}[after-item-skip=6pt](2)
  \task $3x\cdot 4x^{2}-5x^{3}$
  \task $2a^{2}\cdot (-3a)+10a^{4}:(2a)$
  \task $(2x)^{3}-x\cdot 6x^{2}$
  \task $15x^{4}:(3x^{2})-(2x)^{2}$
  \task $-x^{2}\cdot x^{3}+(x^{2})^{2}\cdot 3x$
  \task $(-3x)^{2}\cdot x-4x^{3}$
\end{tasks}

\ex{3} Opera y reduce:\begin{tasks}[after-item-skip=6pt](2)
  \task $(-2x^{2}y)^{3}:(4x^{4}y^{2})$
  \task $(3a^{2}b)^{2}\cdot (-ab)^{3}:(9a^{5}b^{4})$
  \task $(-x^{2})^{3}+2x^{4}\cdot (-x)^{2}-(2x^{3})^{2}:4$
  \task $(2x)^{4}:(-x)^{2}-3x\cdot (-2x)$
\end{tasks}

\ex{3} Un ortoedro tiene las aristas $2x$, $3x$ y $x^2$ cm. Escribe como monomios o polinomios: a) el volumen; b) el área total. c) ¿Qué le pasa al volumen si todas las aristas se hacen el doble de largas?


% =====================================================
\seccio{4}{Polinomios: partes}

\ex{1} Completa la tabla:\par\smallskip{\centering\renewcommand{\arraystretch}{1.4}\small\begin{tabular}{|c|C{3.8cm}|C{1cm}|C{1.7cm}|C{1.7cm}|C{1.5cm}|}\hline Polinomio & Ordenado & Grado & Coef. principal & Término indep. & Términos\\\hline $3x^{2}-5x+1$ & & & & & \\\hline $7-2x^{4}+x$ & & & & & \\\hline $x^{3}-8$ & & & & & \\\hline $-x^{5}+4x^{3}-x^{2}+2x$ & & & & & \\\hline $6x$ & & & & & \\\hline $4x^{2}-x+3x^{3}-2$ & & & & & \\\hline\end{tabular}\par}\recorda{el \textbf{grado} de un polinomio es el grado más alto de sus términos; el \textbf{coeficiente principal}, el del término de mayor grado; el \textbf{término independiente}, el que no tiene letra (con su signo).}

\ex{1} Reduce y ordena estos polinomios. ¿Qué grado tiene cada uno?\begin{tasks}[after-item-skip=6pt](2)
  \task $2x^{2}+3x-x^{2}+5-x$
  \task $x^{3}-4x+2x^{3}+x^{2}-3x^{3}+4x$
  \task $5-x+3x^{2}-7+x-x^{2}$
  \task $-x^{4}+2x+x^{4}-3x^{2}+x^{3}$
\end{tasks}

\ex{1} ¿Cuáles de estos polinomios son completos? Escribe los que no lo son poniendo ceros en los términos que faltan.\begin{tasks}(2)
  \task $x^{3}+2x^{2}-x+5$
  \task $4x^{2}-1$
  \task $x^{4}-x^{3}+x^{2}-x+1$
  \task $2x^{3}+x$
\end{tasks}

\ex{2} Calcula el valor numérico de $P(x)=x^2-3x+2$ y de $Q(x)=2x^3-x^2+4$ para los valores indicados:\begin{tasks}[after-item-skip=5pt](4)
  \task $P(0)$
  \task $P(1)$
  \task $P(-1)$
  \task $P(2)$
  \task $Q(0)$
  \task $Q(1)$
  \task $Q(-2)$
  \task $Q\!\left(\tfrac{1}{2}\right)$
\end{tasks}

\ex{2} Escribe un polinomio que cumpla cada condición: a) de grado 3 y con tres términos; b) de grado 2, completo y con término independiente $-4$; c) de grado 4, con coeficiente principal $-2$ y sin término independiente.

\ex{3} a) Halla $k$ para que $P(x)=x^3-2x^2+kx+6$ cumpla $P(2)=0$. \ b) Halla $a$ si $Q(x)=ax^2+3x-1$ y $Q(2)=13$.

\ex{3} Un polinomio de grado 2 tiene el término independiente igual a 5, el coeficiente de $x$ igual a $-4$ y cumple $P(1)=3$. ¿Cuál es? ¿Cuánto vale $P(-1)$?


% =====================================================
\seccio{5}{Suma y resta de polinomios}

\ex{1} Dados $P(x)=3x^2-5x+2$, \ \ $Q(x)=x^3-2x^2+4x-1$ \ \ y \ \ $R(x)=-x^2+x-6$, calcula:\begin{tasks}(4)
  \task $P+Q$
  \task $P+R$
  \task $Q+R$
  \task $P+Q+R$
\end{tasks}

\ex{1} Con los mismos polinomios, calcula:\begin{tasks}(4)
  \task $P-Q$
  \task $Q-P$
  \task $P-R$
  \task $R-Q$
\end{tasks}\recorda{para restar, cambiamos el signo de \textbf{todos} los términos del polinomio que restamos: $-(x^3-2x^2+4x-1)=-x^3+2x^2-4x+1$.}

\ex{1} Escribe el opuesto de cada polinomio y comprueba que su suma es 0:\begin{tasks}(3)
  \task $2x^{2}-7x+1$
  \task $-x^{3}+x-5$
  \task $4x^{4}-3x^{2}$
\end{tasks}

\ex{2} Con los polinomios del ejercicio 33, calcula:\begin{tasks}(4)
  \task $P+Q-R$
  \task $P-(Q+R)$
  \task $Q-P-R$
  \task $2P-R$
\end{tasks}

\ex{2} Opera y reduce:\begin{tasks}[after-item-skip=6pt](1)
  \task $(x^{2}-4x+3)-(2x^{2}+x-5)+(x^{2}-1)$
  \task $(5x^{3}-2x)-(x^{3}+x^{2}-2x+4)$
  \task $-(x^{2}-3)+(4x^{2}-x)-(2x+7)$
  \task $(2x^{2}y-xy+3)-(x^{2}y+4xy-1)$
\end{tasks}

\ex{3} Halla el polinomio $X(x)$ en cada caso: \ a) $X(x)+(2x^2-3x+1)=5x^2+x-4$; \ b) $(x^3-x+2)-X(x)=3x^2+x$.

\ex{3} Los lados de un pentágono miden $x^2+1$, $2x$, $x^2-x+3$, $3x+2$ y $2x^2-4$ cm. Escribe su perímetro y calcúlalo para $x=4$. ¿Cuánto más largo es el lado más largo que el más corto cuando $x=4$?


% =====================================================
\seccio{6}{Producto de polinomios}

\ex{1} Multiplica:\begin{tasks}[after-item-skip=6pt](3)
  \task $3(x^{2}-2x+5)$
  \task $-2(4x^{3}-x+1)$
  \task $5(x^{2}-3x)$
  \task $-(2x^{2}-x-7)$
  \task $\tfrac{1}{2}(4x^{2}-6x+2)$
  \task $-4(-x^{3}+2x^{2}-3)$
\end{tasks}\recorda{propiedad distributiva: multiplicamos el número o el monomio por \textbf{cada} término del polinomio.}

\ex{1} Multiplica:\begin{tasks}[after-item-skip=6pt](3)
  \task $2x(x^{2}-3x+4)$
  \task $-3x^{2}(2x-5)$
  \task $x^{3}(x^{2}+x-1)$
  \task $-4x(x^{3}-2x+3)$
  \task $5xy(x-2y+1)$
  \task $-x^{2}(-x^{2}+3x-6)$
\end{tasks}

\ex{1} Multiplica y reduce:\begin{tasks}[after-item-skip=6pt](3)
  \task $(x+2)(x+5)$
  \task $(x-3)(x+4)$
  \task $(x-1)(x-6)$
  \task $(2x+1)(x-3)$
  \task $(3x-2)(2x+5)$
  \task $(x^{2}+1)(x-2)$
\end{tasks}\recorda{cada término del primero por cada término del segundo: $(x+3)(2x-5)=2x^2-5x+6x-15=2x^2+x-15$.}

\ex{2} Multiplica y reduce:\begin{tasks}[after-item-skip=6pt](2)
  \task $(x+2)(x^{2}-3x+1)$
  \task $(2x-1)(x^{2}+4x-3)$
  \task $(x^{2}-x+2)(3x+1)$
  \task $(x^{2}+2x-1)(x^{2}-2)$
\end{tasks}

\ex{2} Opera y reduce:\begin{tasks}[after-item-skip=6pt](2)
  \task $3(x+2)-2(x-1)$
  \task $x(x+4)-2(x^{2}-3x)$
  \task $(x+1)(x-2)+(x+3)(x-1)$
  \task $2x(x-5)-(x+2)(x-3)$
\end{tasks}

\ex{2} Escribe, como polinomio reducido, el área de: a) un rectángulo de lados $x+4$ y $2x-1$; b) un cuadrado de lado $x+3$ al que quitamos un cuadrado de lado $x$; c) un triángulo de base $2x+6$ y altura $x$.

\ex{3} Multiplica (haz primero el producto de dos factores y después multiplica por el tercero):\begin{tasks}[after-item-skip=6pt](2)
  \task $(x+1)(x-1)(x+3)$
  \task $(x-2)(2x+1)(x+2)$
\end{tasks}

\ex{3} Opera y reduce:\begin{tasks}[after-item-skip=6pt](1)
  \task $(2x-1)(x+3)-(x-2)(x^{2}+x)+x^{3}$
  \task $(x^{2}-2x+3)(x+1)-x(x^{2}-x+1)$
\end{tasks}


% =====================================================
\seccio{7}{Identidades notables}

\ex{1} Desarrolla (cuadrado de una suma):\begin{tasks}[after-item-skip=6pt](4)
  \task $(x+3)^{2}$
  \task $(x+7)^{2}$
  \task $(2x+1)^{2}$
  \task $(a+5b)^{2}$
\end{tasks}\recorda{$(a+b)^2=a^2+2ab+b^2$: \ el cuadrado del primero, \textbf{más el doble del primero por el segundo}, más el cuadrado del segundo.}

\ex{1} Desarrolla (cuadrado de una diferencia):\begin{tasks}[after-item-skip=6pt](4)
  \task $(x-4)^{2}$
  \task $(x-1)^{2}$
  \task $(3x-2)^{2}$
  \task $(5-x)^{2}$
\end{tasks}\recorda{$(a-b)^2=a^2-2ab+b^2$.}

\ex{1} Desarrolla (suma por diferencia):\begin{tasks}[after-item-skip=6pt](4)
  \task $(x+5)(x-5)$
  \task $(x-9)(x+9)$
  \task $(2x+3)(2x-3)$
  \task $(4+x)(4-x)$
\end{tasks}\recorda{$(a+b)(a-b)=a^2-b^2$: suma por diferencia es igual a diferencia de cuadrados.}

\ex{2} Desarrolla usando las identidades notables:\begin{tasks}[after-item-skip=6pt](3)
  \task $(3x+4)^{2}$
  \task $(2x-5y)^{2}$
  \task $(x^{2}+3)^{2}$
  \task $(x^{2}-2x)^{2}$
  \task $(3a+2b)(3a-2b)$
  \task $(\tfrac{1}{2}x-2)^{2}$
\end{tasks}

\ex{2} Escribe como cuadrado de un binomio o como producto de una suma por una diferencia:\begin{tasks}[after-item-skip=6pt](3)
  \task $x^{2}+10x+25$
  \task $x^{2}-6x+9$
  \task $x^{2}-16$
  \task $4x^{2}+4x+1$
  \task $9x^{2}-25$
  \task $x^{2}-2x+1$
\end{tasks}

\ex{2} Completa las casillas:\begin{tasks}[after-item-skip=6pt](2)
  \task $(x+\square)^2=x^2+\square+16$
  \task $(x-\square)^2=x^2-10x+\square$
  \task $(\square+3)^2=4x^2+\square+9$
  \task $(x+\square)(x-\square)=x^2-49$
\end{tasks}

\ex{3} Simplifica:\begin{tasks}[after-item-skip=6pt](1)
  \task $(x+2)^{2}-(x-2)^{2}$
  \task $(2x-1)^{2}-(2x+1)(2x-1)$
  \task $(x+3)^{2}-2(x-1)(x+1)-x(x+6)$
\end{tasks}

\ex{3} Calcula sin calculadora, usando las identidades notables:\begin{tasks}(4)
  \task $41^2$
  \task $29^2$
  \task $103\cdot97$
  \task $55^2-45^2$
\end{tasks}

\ex{3} a) Demuestra que $(a+b)^2-(a-b)^2=4ab$. \ b) Úsala para calcular $(x+5)^2-(x-5)^2$ y $\ 1001^2-999^2$.


% =====================================================
\seccio{8}{Extracción de factor común}

\ex{1} Saca factor común:\begin{tasks}[after-item-skip=6pt](3)
  \task $3x+12$
  \task $5x^{2}-10x$
  \task $x^{3}+x^{2}$
  \task $8x-12$
  \task $7x^{2}+14x^{3}$
  \task $6x^{3}-4x^{2}+2x$
\end{tasks}\recorda{el factor común es el MCD de los coeficientes por las letras comunes a todos los términos con el \textbf{menor} exponente. Dentro del paréntesis va lo que queda al dividir cada término entre el factor común (si un término es igual al factor común, queda un 1).}

\ex{1} Calcula de la manera más rápida, sacando factor común:\begin{tasks}(4)
  \task $34\cdot 18+34\cdot 82$
  \task $125\cdot 7-125\cdot 3$
  \task $9\cdot 47+9\cdot 53$
  \task $56^2-56\cdot 46$
\end{tasks}

\ex{2} Saca factor común:\begin{tasks}[after-item-skip=6pt](2)
  \task $4a^{2}b-6ab^{2}$
  \task $9x^{4}-6x^{3}+3x^{2}$
  \task $10x^{3}y^{2}+15x^{2}y^{3}$
  \task $x^{5}-x^{4}+x^{3}-x^{2}$
  \task $-8x^{3}-12x$
  \task $12a^{2}b^{3}-18a^{3}b^{2}+6a^{2}b^{2}$
\end{tasks}

\ex{2} Completa las casillas:\begin{tasks}[after-item-skip=6pt](2)
  \task $6x^3-9x^2=3x^2(\square-\square)$
  \task $10x^4+5x^2=5x^2(\square+\square)$
  \task $x^2y-xy^2=\square(x-y)$
  \task $14a^3-21a^2b=\square(2a-3b)$
\end{tasks}

\ex{3} Factoriza tanto como puedas (primero el factor común y después una identidad notable):\begin{tasks}[after-item-skip=6pt](2)
  \task $2x^{3}-8x$
  \task $3x^{3}+6x^{2}+3x$
  \task $x^{4}-9x^{2}$
  \task $5x^{3}-20x^{2}+20x$
\end{tasks}

\ex{3} Saca factor común (el factor común es un binomio):\begin{tasks}[after-item-skip=6pt](3)
  \task $x(x+2)+3(x+2)$
  \task $2x(x-5)-(x-5)$
  \task $(x+1)^{2}+4(x+1)$
\end{tasks}

\ex{3} Sin calculadora: a) $99^2+99$; \ b) $13\cdot87+13^2$; \ c) Demuestra que $n^2+n$ siempre es par para cualquier número natural $n$.

\vfill
\begin{center}\color{muted}\small\faCheckCircle[regular]\ \ ¡Ánimo! Comprueba tus resultados con las herramientas interactivas de los apuntes de la unidad.\end{center}

\end{document}
